The first display is immediate from the matched lower and upper bounds in \(\text{thm:adaptive-rootn-minimax}\), whose constants are uniform in \(n\).  The condition \(\varepsilon_{1,n}\le\varepsilon_0\) permits a fixed positive sequence.  Inspection of \(\text{def:non-gaussian-class}\) and the effective construction in \(\text{def:adaptive-contour-estimator}\) shows that only \(\bar g_n\), and neither an outcome code nor \(\varepsilon_{2,n}\), enters the spectral branch.

For the separate Gaussian intersection, \(\text{prop:bounded-outcome-gaussian-degeneracy}\) gives \(\theta_0(P)=0\) for every member.  If the class is nonempty, the zero statistic has zero squared and generalized-quantile loss.  Nonnegativity and \(\text{def:minimax-risks}\) therefore give both zero minimax criteria.  The definitions assert no equality or inclusion between the two classes.  The Gaussian conclusion is solely a diagnostic of the source theorem's simultaneous assumptions, and the shrinking-separation and inference questions remain untouched.